
\uline{un nombre premier} $p$ \uline{et} $G''$ \uline{est un groupe de torsion. Alors toute application polynôme} $f:G\to\mathbb Q$ \uline{qui se factorise par} $\mathbb Z$\uline{, est constante.}

\emph{Démonstration} : Il s'agit de montrer que pour tout $x\in G$, on a $f(x)=f(0)$.

a) Supposons d'abord que $x$ soit un élément de torsion, disons $mx=0$ ($m\in Z$, $m\ne0$). Alors le polynôme $P_x(n)=f(nx)-f(0)$ a un nombre infini de racines, à savoir tous les entiers divisibles par $m$. Donc le polynôme $P_x$ est identiquement nul et on a $f(x)=f(0)$.

b) Supposons maintenant que $x$ soit un élément libre de $G$ divisible par toutes les puissances $p^k$ ($k\in\mathbb Z$). Considérons des polynômes $P_k$ définis par
\[
P_k(n)=f(nx/p^k)=\sum_{i\in\mathbb N}n^iy(x/p^k;i)\in\mathbb Z
\qquad(n\in\mathbb Z;\ y(x/p^k,i)\in\mathbb Q).
\]
Pour tout $n\in\mathbb Z$ on a l'identité
\[
P_k(n)=P_{k+1}(pn)
\]
i.e.
\[
\sum_{i\in\mathbb N}n^iy(x/p^k;i)=\sum_{i\in\mathbb N}p^in^iy(x/p^{k+1};i),
\]
d'où on tire par comparaison de coefficients
\[
y(x/p^k;i)=p^iy(x/p^{k+1};i)\qquad(i\in\mathbb N).
\]
